\begin{teilaufgaben}
\item
The domain $\Omega$ is the rectangle $[0,2] \times [0, 3]$. 
The function $u(x,y)$ is defined on $\Omega$.
It satisfies in $\Omega$ 
\[
\Delta u(x,y) = 0.
\]
On the boundary of $\Omega$ one has $u(x,y) = 2 \cdot x - x^2$. 

Determine approximate values for $u(1,1)$ and $u(1,2)$. 
Use finite volumes \`a la Voronoi.  

\item
b) The boundary of the domain $\Omega$ is a polygon with vertices
\[
(0,0), (4,0), 
(4, 2), (2,2), (2,17), (0,17), (0,0).
\]
The function $u(x,y)$ is defined on $\Omega$. It satisfies in $\Omega$ 
\[
\Delta u(x,y) = 1.
\]
On the boundary of $\Omega$ one has $u(x,y) = 0$.

Determine approximate values for $u(1,1)$, $u(3,1)$ and $u(1,3)$.
Use finite volumes with three Voronoi cells.   
\end{teilaufgaben}


\begin{loesung}
\begin{teilaufgaben}
\item
By the finite volume method we have
\[
3/2 \cdot \big(0 - \tilde u(1,1)\big)
+
2 \cdot \big(\tilde u(1,2) - \tilde u(1,1)\big)
+
3/2 \cdot \big( 0 - \tilde u(1,1)) +  2 \cdot (1 - \tilde u(1,1)\big)
=
0
\]
or
\begin{equation}
-14 \cdot \tilde u(1,1) + 4 \cdot \tilde u(1,2) = -4.
\tag{\bf{1P}}
\end{equation}
By symmetry we have
\begin{equation}
-14 \cdot \tilde u(1,2) + 4 \cdot \tilde u(1,1) = -4.
\tag{\bf{1P}}
\end{equation}
Therefore
\begin{equation}
\big(\tilde u(1,1), \tilde u(1,2)\big) = (0.4, 0.4).
\tag{\bf{1P}}
\end{equation}
\item
By the finite volume method we have
\[
\renewcommand{\arraycolsep}{1pt}
\renewcommand{\arraystretch}{1.4}
\hspace*{-10pt}
\begin{array}{rclcrclcrclcrclcl}
\frac{2}{2}&\cdot&(\tilde u(2,1) - \tilde u(1,1))
&+&
\frac{2}{2}&\cdot&(\tilde u(1,2)  - \tilde u(-1,1))
&+&
\frac{2}{1}&\cdot&(0 - \tilde u(1,1))
&+&
\frac{2}{1}&\cdot&(0 - \tilde u(1,1))
&=&
4,
\\
\frac{2}{1}&\cdot&(\tilde 0 - \tilde u(2,1))
&+&
\frac{2}{1}&\cdot&(0  - \tilde u(2,1))
&+&
\frac{2}{2}&\cdot&(\tilde u(1,1)- \tilde u(2,1))
&+&
\frac{2}{1}&\cdot&(0 - \tilde u(2,1))
&=&
4,
\\
\frac{15}{1}&\cdot&(\tilde 0 - \tilde u(1,2))
&+&
\frac{2}{14}&\cdot&(0  - \tilde u(1,2))
&+&
\frac{15}{1}&\cdot&(0 - \tilde u(1,2))
&+&
\frac{2}{2}&\cdot&(\tilde u(1,1)- \tilde u(1,2))
&=&
30.
\end{array}
\]
This leads to a system of three linear equations:
\begin{equation}
\begin{pmatrix*}[r] -6 & 1 & 1 \\ 1 & -7 & 0 \\  0 & 1 &  -218/7 \end{pmatrix*}
\cdot
\begin{pmatrix} \tilde u(1,1) \\ \tilde u(2,1) \\ \tilde u(1,2) \end{pmatrix}
=
\begin{pmatrix*}[r] 4  \\  4 \\  30 \end{pmatrix*}.
\tag{\bf{2P}}
\end{equation}
Hence
\begin{equation}
\tilde u(1,1) = -0.9502,\qquad \tilde u(2,1) = -0.7072
\qquad\text{and}\qquad
\tilde u(1,2) = -0.9938.
\tag{\bf{1P}}
\end{equation}
\end{teilaufgaben}
\end{loesung}
